% Source PDF: source/普通高考/2013/2013安徽文.pdf, page 1; vector program flowchart.
% Grid evidence: tmp/review_2013_anhui_liberal/grid/page-1-grid100.jpg and
%   q3_original_clean.png (no-grid crop from the source page).
% Elements: rounded start/end terminals; initialization rectangle s=0,n=2;
% decision diamond n<8?; update rectangles s=s+1/n and n=n+2;
% output parallelogram; directed connectors and 是/否 labels.
% Complete node -> directed-edge list from the source:
%   start -> init -> test; test --是--> update_s -> increment -> left loop
%   -> main entry above test; test --否--> output -> finish. The loop return
% has one right-pointing arrow into the main vertical stem and remains separate
% from the terminating branch; all borders/connectors are solid.

\begin{tikzpicture}[
  >=Stealth,
  line width=0.45pt,
  line join=round,
  line cap=round,
  font=\normalsize,
  flowtext/.style={anchor=center, align=center, inner sep=0pt,
    execute at begin node={\setlength{\parindent}{0pt}}},
  every node/.style={flowtext},
  flowbox/.style={flowtext, draw, minimum height=0.68cm,
    inner xsep=3pt, inner ysep=2pt},
  terminal/.style={flowbox, rounded corners=0.18cm, minimum width=1.45cm},
  process/.style={flowbox},
  decision/.style={flowbox, diamond, aspect=2.7, minimum width=3.50cm,
    minimum height=0.95cm},
  io/.style={flowbox, trapezium, trapezium left angle=72,
    trapezium right angle=72, minimum width=2.35cm, minimum height=0.76cm},
]
  \node[terminal] (start) at (0,8.45) {开始};
  \node[process, minimum width=2.65cm] (init) at (0,6.95) {$s=0,n=2$};
  \node[decision] (test) at (0,5.10) {$n<8?$};
  \node[process, minimum width=3.05cm, minimum height=0.95cm]
    (update_s) at (0,3.40) {$s=s+\dfrac{1}{n}$};
  \node[process, minimum width=2.65cm] (increment) at (0,1.75)
    {$n=n+2$};
  \node[io] (output) at (3.30,3.35) {输出 $s$};
  \node[terminal] (finish) at (3.30,1.80) {结束};

  % Main trunk and terminating branch.
  \draw[->] (start.south) -- (init.north);
  \draw[->] (init.south) -- (test.north);
  \draw[->] (test.south) -- node[flowtext, right] {是} (update_s.north);
  \draw[->] (update_s.south) -- (increment.north);
  \draw[->] (test.east) -- node[flowtext, above] {否}
    (3.30,5.10) -- (3.30,3.78) -- (output.north);
  \draw[->] (output.south) -- (finish.north);

  % The 是 branch loops left and returns with a single horizontal
  % right-pointing arrow into the main stem above the decision diamond.
  \coordinate (loopjoin) at (0,5.95);
  \draw[->] (increment.south) -- (0,0.55) -- (-2.55,0.55)
    -- (-2.55,5.95) -- (loopjoin);
\end{tikzpicture}
